%==============================================================================
\section{Reproducing the experiments}\label{app:repro}
%==============================================================================

The complete implementation --- the solvers, the reimplemented baselines, the campaign
harness, the verification program, and the per-instance result files behind every
number in Section~\ref{sec:comp} --- is available at
\url{https://github.com/shankar-n/JGP-SSP-Codes-Shankar}.

\subsection{Layout}\label{ssec:repro-layout}

\noindent\begin{tabular}{@{}l>{\raggedright\arraybackslash}p{9.4cm}@{}}
\toprule
Path & Contents \\
\midrule
\texttt{src/SSP/}     & the loading oracle, instance handling, heuristics, the grouping solver \\
\texttt{src/BBC/}     & the branch-and-Benders-cut solver, the three reimplemented baselines, and the campaign harness \\
\texttt{src/BNP/}     & the PCF$'$ and PTF branch-and-price prototypes and their pricers \\
\texttt{data/}        & the instance families of Table~\ref{tab:families} \\
\texttt{verification/}& the independent verification program of Appendix~\ref{app:verify} \\
\texttt{report/}      & the source of this document \\
\bottomrule
\end{tabular}

\subsection{Requirements}\label{ssec:repro-req}

The exact solvers require IBM CPLEX with its Python interface; version~22.1.1 was used
throughout. The branch-and-price prototypes require PySCIPOpt. The grouping solver
requires PySCIPOpt. The verification program of Appendix~\ref{app:verify} requires
only Python~3 and SciPy, and in particular needs no commercial solver, so the
structural results of Section~\ref{sec:gap} can be checked on any machine.

\subsection{Instance format}\label{ssec:repro-format}

An instance file gives the number of jobs, the number of tools and the capacity on one
line, followed by the job--tool incidence matrix with one row per tool and one column
per job, where entry $(t,j)$ is $1$ when job $j$ requires tool $t$.

\subsection{Running the campaign}\label{ssec:repro-campaign}

The instance families and the solver configurations are declared in one place, in the
campaign configuration module, which is the single source of truth for both. The runner
is resumable: it records one row per completed run and skips runs already recorded. Its
completion key is (family, file name, jobs, tools, capacity, configuration), because
file name and configuration do not identify a run in collections with repeated names.
The historical campaign runner used that shorter key, which caused the $58$ skipped
pairs disclosed in Section~\ref{ssec:bench}; the repository version and its regression
check contain the correction, and same-protocol completion filled all $58$. If the solver changes, previously recorded rows must also
be deleted before re-running, or the runner will keep stale results. Both re-runs
described in this report --- the one following the Benders cut correction and the one
following the addition of the coverage row --- required exactly that.

On the cluster the campaign is submitted as a SLURM job array with one task per
configuration and instance shard. Each run is given four dedicated threads,
$16$ gigabytes of memory and a one-hour limit.

\subsection{Reproducing the numbers in Section~\ref{sec:comp}}\label{ssec:repro-numbers}

Every table in Section~\ref{sec:comp} is computed from the raw per-instance result
files by a single analysis pass, which reads the shards, keys each row by the
five-part identifier (family, file name, jobs, tools, capacity), asserts the raw total
of $17{,}052$ rows, and excludes the $132$ rows attached to the eleven non-canonical
Catanzaro files identified in Section~\ref{ssec:bench}. It then asserts the canonical
total of $1{,}410$ instances and $16{,}920$ rows before deriving:

\begin{itemize}[topsep=3pt,itemsep=1pt]
  \item the solved counts of Table~\ref{tab:solves}, per family and in total, against
  the fixed $1{,}410$ planned identities per configuration;
  \item the cross-method agreement check, comparing the empty-start cost of the
  returned sequence across every pair of methods that both solved an instance;
  \item the shifted geometric mean times of Table~\ref{tab:times}, on the subset that
  all compared configurations solve;
  \item the fractional-cut comparison of Table~\ref{tab:frac} and the ablation of
  Table~\ref{tab:ablation}, each on the instances that both compared configurations
  attempted;
  \item the coverage-bound split of Table~\ref{tab:tightness}, by reading each instance
  file, counting the tools $\lvert\U\rvert$ that some job actually requires, and testing
  whether the known optimum equals $\lvert\U\rvert$ in the empty-start cost;
  \item the branch-and-price table and root-bound counts of
  Section~\ref{ssec:results-bnp}, from the separate result files of the prototypes.
\end{itemize}

The figures are generated by \texttt{report/make\_figure\_data.py}, which writes the
five \texttt{.dat} files that Sections~\ref{sec:intro}, \ref{sec:comp} and
\ref{sec:disc} plot. Three of the figures read those files directly. Two ---
Figure~\ref{fig:tightsplit} and Figure~\ref{fig:rootgap} --- carry their coordinates
inline in the \LaTeX{} instead, because their $x$ axes are categorical. Their values
were checked against \texttt{solvedvsgap.dat} and \texttt{rootgap.dat} point by point,
but they will not follow a regeneration on their own and have to be updated by hand if
the campaign is re-run.

Two points of the key matter, because getting either wrong changes the
conclusions. The capacity must be part of the identifier: the Crama collection reuses
file names across four capacities, and keying on the name alone merges rows from
different instances and manufactures cross-method disagreements that do not exist. And
the coverage bound must be computed from $\lvert\U\rvert$ rather than from the tool
count declared in the instance file, because ten instances of the collection declare a
tool that no job requires.

Headline solve comparisons use the fixed $1{,}410$-instance denominator, so every
non-optimal outcome counts as unsolved. For reproducibility, the primary comparison uses
the resource envelope stated above. An auxiliary LSS sensitivity
used 32~GiB; it certifies no additional optimum and is never merged with the primary
ledger. The public campaign analysis in
\path{verification/analyse_campaign_results.py} validates all $16{,}920$ canonical
identities, while
\path{src/BBC/results/provenance/primary_ledger_audit.json} provides a compact SHA-256
inventory of the pre-merge primary archive. Together they make the denominator and its
provenance auditable without publishing the private cluster workspace.
